% ── §7b Methods: Conservation Analysis ──
\section{Conservation Analysis}
\label{sec:methods-conservation}

\subsection{Column-Level Conservation}
For each alignment column $c$, we compute the Shannon entropy $H(c)$ as
described in Section~\ref{sec:theory-info}.  A column is classified as
\emph{conserved} if $H(c) < 0.3$ bits, following the convention used
throughout our pipeline.  This threshold captures positions where a single
amino acid accounts for more than approximately 90\% of observations.

\subsection{Pair-Level Conservation Types}
Each candidate pair $(i, j)$ is assigned one of four conservation types based
on the binary conservation status of its two constituent columns:

\begin{center}
	\begin{tabular}{clc}
		\toprule
		Type      & Description                                   & Encoding \\
		\midrule
		var/var   & both columns variable                         & 00       \\
		var/cons  & position $i$ variable, position $j$ conserved & 01       \\
		cons/var  & position $i$ conserved, position $j$ variable & 10       \\
		cons/cons & both columns conserved                        & 11       \\
		\bottomrule
	\end{tabular}
\end{center}

These two bits serve as additional feature dimensions in the complete circuit,
expanding the feature space from 8 to 10 bits (1024 cells).

\subsection{Soft Conservation Gradient}
Rather than using only the binary threshold, we also analyse contact rates as
a continuous function of the minimum column entropy within each pair:
$\min(H(i), H(j))$.  This analysis reveals a non-monotone relationship:
fully conserved pairs (entropy $\approx 0$) show \emph{below-random} contact
rates ($0.69\times$), while near-conserved pairs (entropy $0.001$--$0.05$)
show above-random enrichment ($1.35\times$).  The most variable quintile shows
the lowest enrichment ($0.70\times$).  The overall Spearman correlation between
column entropy and contact status is negative ($\rho = -0.053$, $p \approx 0$),
confirming that conservation carries structural information.

The non-monotonicity at the extreme conserved end likely reflects the presence
of catalytic or functional residues (such as active-site serines or
disulfide-forming cysteines) that are absolutely conserved for functional
rather than structural reasons and may project away from the structural core.

\subsection{Cross-Protein Transfer of Negative Selection}
A key prediction of our framework is that cells marked as FORBIDDEN by the
flipped Karnaugh map---combinations never observed across training
proteins---should remain rare on held-out proteins if the underlying
selection pressure is universal rather than family-specific.  We test this by
measuring the contact rate within forbidden-called cells on test proteins and
comparing it to the training-set prediction.

Across five folds of protein-level cross-validation, the observed contact rate
within forbidden-called cells on test proteins averages $1.7\%$, matching the
training-set upper bound of $1.75\%$ (= base$/2$) almost exactly.  This
represents a specificity of $98.3\%$ and demonstrates that negative-selection
rules generalise across protein families with high fidelity.
